% TT09032C
% Rückenmarksverletzung mit Symptomen
% 14.0
% 2013-11-30
\label{m:ws-trauma}
\EE{Taktik}: \Speech{\Ca{Zeitkritisch! Vitale Bedrohung!} Zügige Bergung,
Immobilisation und Transport in eine geeignete Einrichtung}


Vgl. \prettyref{m:einschaetzung-indikation-ws-immobilisation}
\begin{itemize}
  \item Bei ansprechbaren Patienten: immer
  \EE{\Abk{HWS}-Immobilisation} (HWS-Schiene)
  \item \EE{Bergung} mit Schaufeltrage, Spineboard oder
  KED-System
  \item \EE{Lagerung} und Schienung mittels
  mit Vakuummatratze, Schaufeltrage mit Fixationsmöglichkeit oder
  Spineboard.\footnote{Welches Hilfsmittel
  eingesetzt wird ist abhängig vom Patientenzustand,
  dem Training des Teams, internen Vorschriften und
  vorhandenem Material.}
  \item Bei Querschnittslähmung: Kennzeichnung der
  Sensibilitätsgrenze mit Angabe der Uhrzeit
  \item \EE{Monitoring}: Achte besonders auf Störungen
  von Atmung und Bewußtsein, sowie Schockentwicklung.
%   \item \EE{Unfallmechanismus} erheben und dokumentieren 
  \item \EE{Transportentscheidung}: Schockraum
\end{itemize}